\chapter{Financiamiento, Dilución y \emph{Runway}}\label{ch:fundraising-and-dilution}

% Alcance: \emph{cap tables} y dilución (pre/post-money), SAFEs y notas convertibles,
%   \emph{term sheets} y preferencias de liquidación, métricas para inversores y el
%   múltiplo de ARR. Capítulo nuevo (cierra la brecha entre el fundador y las finanzas).
%   Sigue a valoración (\cref{ch:valuation}) y estructura de capital
%   (\cref{ch:risk-and-capital-structure}): del valor empresarial a cómo los fundadores
%   financian en la práctica.
% Referencias cruzadas: eq:scenario-ev (ancla del caso pesimista del Seed pre-money, ch.20);
%   eq:dcf / eq:ev-revenue (cadena de valoración reutilizada como marco del inversor);
%   eq:rule40 / eq:burn-multiple / eq:runway (métricas de eficiencia de los ch.21/26
%   que impulsan la narrativa de financiamiento).

\section{Cap Tables y Dilución de Propiedad}\label{sec:cap-table}
% complexity: 3

¿Qué le cuesta al fundador levantar una ronda? La \emph{cap table} --- el registro de
cada acción, opción e instrumento convertible --- responde a esa pregunta con precisión, y
su aritmética no perdona. Entender la dilución antes de firmar significa dominar la fórmula
post-money de memoria.

Cuando una empresa emite nuevas acciones para inversores a una valoración acordada, el
porcentaje de los dueños existentes disminuye aunque el valor total de su participación
pueda aumentar. Las dos cantidades que un fundador debe rastrear son (i) la \emph{valoración
post-money}, que fija el precio por acción, y (ii) el \emph{porcentaje de propiedad del
fundador} tras absorber las nuevas acciones en el conteo totalmente diluido. La identidad
central es:
\begin{equation}\label{eq:post-money}
\begin{gathered}
  \mathrm{PostMoney} \;=\; \mathrm{PreMoney} \;+\; \mathrm{Inversión}, \\
  \text{Inversor\%} \;=\; \frac{\mathrm{Inversión}}{\mathrm{PostMoney}}, \qquad
  \text{Fundador\%}_{\mathrm{post}} \;=\; \text{Fundador\%}_{\mathrm{pre}}\,\bigl(1 - \text{Inversor\%}\bigr).
\end{gathered}
\end{equation}
Los fundadores comienzan en \qty{100}{\percent} y se diluyen de forma multiplicativa en
cada ronda. Cada financiamiento sucesivo multiplica el porcentaje superviviente por
\((1 - \text{nuevo inversor\%})\); los porcentajes de todos los dueños existentes ---
fundadores, ángeles e inversores de rondas previas --- se reducen proporcionalmente para
dar cabida al capital entrante.

Una sutileza que la \cref{eq:post-money} oculta es el \emph{option pool shuffle}. Los
inversores suelen exigir que el grupo de opciones para empleados no asignadas se amplíe
\emph{antes} de contabilizar el nuevo dinero, lo que significa que la dilución recae
íntegramente sobre los fundadores en lugar de compartirse con el nuevo inversor
\parencite{feld2019venture}. Un \emph{term sheet} que exige un grupo del \qty{15}{\percent}
post-money en una empresa con pre-money de \money{10000000} obliga al fundador a absorber
la dilución del grupo antes de que se fije el precio por acción. Los fundadores que ignoran
esta aritmética la pagan dos veces.

\begin{example}[Cap table de dos rondas]
El valor empresarial ponderado por probabilidad de escenario de la empresa se estimó en
aproximadamente \money{16500000} (véase \cref{eq:scenario-ev}), pero el caso pesimista de
\money{8000000} es el ancla conservadora para la negociación del pre-money Seed. La empresa
levanta una ronda Seed a una valoración pre-money de \money{6000000} con una inversión de
\money{2000000}.

Aplicando \cref{eq:post-money}:
\[
  \mathrm{PostMoney}_{\text{Seed}} \;=\; \money{6000000} + \money{2000000} \;=\; \money{8000000}.
\]
La participación del inversor Seed: \(\money{2000000}/\money{8000000} = \qty{25}{\percent}\).
La participación del fundador cae de \qty{100}{\percent} a \qty{75}{\percent}.

Luego la empresa levanta una \emph{Series~A} a una valoración pre-money de \money{10000000},
emitiendo \money{2000000} de nueva equity.
\[
  \mathrm{PostMoney}_{\text{Series A}} \;=\; \money{10000000} + \money{2000000} \;=\; \money{12000000}.
\]
La participación del inversor Series~A: \(\money{2000000}/\money{12000000} \approx \qty{16.7}{\percent}\).
Aplicando la regla multiplicativa de \cref{eq:post-money}:
\[
  \text{Fundador\%}_{\text{post-A}}
  \;=\; \qty{75}{\percent} \;\times\; (1 - 0.167)
  \;=\; \qty{75}{\percent} \;\times\; 0.833
  \;\approx\; \qty{62.5}{\percent}.
\]
La cap table post-Series~A acumulada: fundador \qty{62.5}{\percent}; inversor Seed
\qty{20.8}{\percent}; inversor Series~A \qty{16.7}{\percent}. Total: \qty{100}{\percent}.

La implicación no es que el fundador «perdió» el \qty{37.5}{\percent}: en un post-money de
\money{12000000}, su \qty{62.5}{\percent} vale \money{7500000} --- más de lo que valía la
empresa entera antes del Seed.
\end{example}

\begin{admonition}{Advertencia}
El porcentaje que retiene un fundador tras una ronda es casi irrelevante en aislamiento. Lo
que importa es el resultado absoluto: el \qty{62.5}{\percent} de una salida de
\money{200000000} vale más que el \qty{90}{\percent} de una salida de \money{10000000} por
un factor de 11. Optimizar la retención de propiedad a expensas del capital que acelera el
crecimiento es un error frecuente del fundador. El marco correcto es el valor esperado de la
posición en dólares del fundador, no la fracción de la torta --- una distinción que la lógica
DCF de \cref{eq:dcf} hace rigurosa.
\end{admonition}

La dilución se acumula entre rondas como el producto de las fracciones supervivientes; los
grupos de opciones, warrants e instrumentos convertibles expanden aún más el denominador.
El cuadro completo requiere rastrear el conteo de acciones \emph{totalmente diluido} en
lugar de sólo las acciones emitidas.
\textcite{metrick2021venture} desarrollan la mecánica completa de múltiples rondas, incluidos
los ajustes por grupo de opciones y warrants, y muestran que los fundadores de etapa tardía
con inversores estratégicos suelen retener entre \qty{10}{\percent} y \qty{30}{\percent} al
momento de la salida a bolsa.

\section{SAFEs y Notas Convertibles}\label{sec:safes-convertibles}
% complexity: 3

Las empresas en etapa temprana raramente tienen el historial operativo para justificar una
negociación dura de pre-money. Dos instrumentos difieren esa conversación: el Simple
Agreement for Future Equity (SAFE), introducido por Y Combinator en 2013, y la nota
convertible, más antigua. Ambos inyectan capital de inmediato y se convierten en equity en
una ronda futura con precio definido, pero sus estructuras económicas difieren materialmente.
Una nota convertible es deuda --- acumula intereses, tiene fecha de vencimiento y puede
forzar al fundador a incumplimiento si no llega una ronda con precio a tiempo. Un SAFE no es
deuda; no acumula intereses ni tiene vencimiento. La protección económica para los
inversores tempranos proviene en cambio de dos parámetros: un \emph{tope de valoración}
(la valoración máxima a la que el SAFE se convierte, independientemente de cuán alto se
fije el precio de la Series~A) y una \emph{tasa de descuento} (una reducción porcentual
sobre el precio por acción de la Series~A). La mayoría de los SAFEs usan uno o ambos.

El precio de conversión --- el precio por acción que paga el inversor del SAFE --- es:
\begin{equation}\label{eq:safe-conversion}
  P_{\mathrm{conv}}
  \;=\; \min\!\bigl(\mathrm{Cap},\;\mathrm{SeriesA}_{\mathrm{pre}}\bigr)
        \;\times\; (1 - \delta),
  \qquad
  N_{\mathrm{new}}
  \;=\; \frac{\mathrm{Inversión}}{P_{\mathrm{conv}}},
\end{equation}
donde \(\delta\) es la tasa de descuento y \(N_{\mathrm{new}}\) es el número de nuevas
acciones emitidas al titular del SAFE al convertirse. Cuando el pre-money de la Series~A
supera el tope, el tope es vinculante; cuando cae por debajo del tope, el precio de la
Series~A (descontado por \(\delta\)) rige. El límite que produce el mayor número de
acciones --- y por lo tanto el menor precio efectivo por acción --- será siempre el que el
inversor prefiera.

\begin{example}[Conversión de SAFE con tope y descuento]
Un ángel invierte \money{500000} en un SAFE con un tope de \money{5000000} y un descuento
del \qty{20}{\percent}. La empresa levanta posteriormente una Series~A a una valoración
pre-money de \money{12000000}.

Primero, determinar qué límite rige:
\begin{itemize}
  \item Valoración efectiva basada en el tope: \money{5000000}.
  \item Valoración efectiva basada en el descuento: \(\money{12000000} \times (1 - 0.20) = \money{9600000}\).
\end{itemize}
El tope (\money{5000000}) es menor, por lo que el tope es vinculante --- el inversor del
SAFE convierte como si el pre-money fuera sólo \money{5000000}. Aplicando
\cref{eq:safe-conversion} con \(\delta = 0\) una vez que el tope rige (el tope ya incorpora
el beneficio económico; usar el tope y el descuento simultáneamente contaría doble):
\[
  P_{\mathrm{conv}} \;=\; \money{5000000} \;\times\; 1.00
                   \;=\; \money{5000000} \text{ (valoración efectiva).}
\]
La fracción de propiedad del titular del SAFE al convertirse:
\[
  \text{SAFE\%}
  \;=\; \frac{\mathrm{Inversión}}{\mathrm{Cap}}
  \;=\; \frac{\money{500000}}{\money{5000000}}
  \;=\; \qty{10}{\percent} \;\text{(base pre-Series-A, post-SAFE).}
\]
Tras el inversor de la Series~A tomar el \qty{16.7}{\percent} (como en el ejemplo anterior),
todos los dueños existentes --- inversor del SAFE y fundadores --- se diluyen
proporcionalmente:
\[
  \text{SAFE\%}_{\text{post-A}}
  \;=\; \qty{10}{\percent} \;\times\; (1 - 0.167)
  \;\approx\; \qty{8.3}{\percent}.
\]
El inversor del SAFE termina con el \qty{8.3}{\percent} de un post-money de
\money{12000000}, un valor implícito de aproximadamente \money{1000000} sobre una inversión
de \money{500000} --- un incremento de \qty{2}{\times} en el momento del cierre de la
Series~A, antes de cualquier salida.
\end{example}

\begin{admonition}{Advertencia}
Múltiples SAFEs con distintos topes y distintos descuentos se convierten a precios
diferentes, emitiendo distintos números de acciones. Los fundadores que levantan un SAFE de
\money{250000} a un tope de \money{3000000}, luego un SAFE de \money{500000} a un tope de
\money{5000000}, y luego un SAFE de \money{1000000} a un tope de \money{8000000} deben
modelar la cascada de conversión completa antes de firmar el \emph{term sheet} de la
Series~A --- la interacción puede ser mucho más dilutiva de lo que cualquier SAFE individual
implicaba. Cada SAFE se resuelve de forma independiente usando \cref{eq:safe-conversion}, y
los conteos de acciones resultantes se acumulan antes de admitir el nuevo dinero
institucional.
\end{admonition}

Los SAFEs y las notas convertibles intercambian economía y velocidad contra complejidad en
el evento de conversión. La economía está gobernada íntegramente por los parámetros del
tope y el descuento; la estructura legal está estandarizada \parencite{feld2019venture}. Una
vez que se fija el precio de la Series~A, el SAFE desaparece en la cap table como acciones
preferentes, sujetas a la misma mecánica de liquidación descrita en la siguiente sección.

\section{Term Sheets y Economía del Inversor}\label{sec:term-sheets}
% complexity: 3

Un \emph{term sheet} de capital de riesgo asigna simultáneamente derechos sobre flujos de
efectivo, derechos de control y protección contra pérdidas. El núcleo económico es la
\emph{preferencia de liquidación}: el derecho de los accionistas preferentes a recuperar su
inversión (o un múltiplo de ella) antes de que los accionistas comunes reciban algo en una
venta, fusión o disolución. Entender la preferencia es entender quién cobra en cada
escenario de salida realista.

Dos estructuras de preferencia dominan. En la \emph{preferencia no participante}, el
inversor elige el mejor de dos resultados: tomar la preferencia establecida o convertir a
acciones comunes y participar pro rata. En la \emph{preferencia participante}, el inversor
toma la preferencia \emph{y luego} participa pro rata en el remanente --- el «doble cobro».
La preferencia no participante es el estándar del mercado; la participante aparece en mercados
competitivos y rondas de dificultad financiera, y es materialmente peor para los fundadores.
La fórmula para los ingresos de salida al inversor es:
\begin{equation}\label{eq:liquidation-preference}
  \mathrm{Inversor}_{\text{no-part}}
  \;=\; \max\!\bigl(\mathrm{Pref},\; p \cdot \mathrm{Salida}\bigr),
  \qquad
  \mathrm{Inversor}_{\text{part}}
  \;=\; \mathrm{Pref} \;+\; p \cdot (\mathrm{Salida} - \mathrm{Pref}),
\end{equation}
donde \(\mathrm{Pref} = m \times \mathrm{Inversión}\) es la preferencia de liquidación
(múltiplo \(m\), normalmente \(1\times\)) y \(p\) es la fracción de propiedad del inversor.
Los fundadores reciben lo que queda una vez satisfechos los inversores.

\begin{example}[Preferencia no participante vs.\ participante]
Un inversor de Series~A invierte \money{5000000} por el \qty{25}{\percent} de la empresa.
La preferencia de liquidación es \(1\times\) no participante. La empresa se vende por
\money{8000000}.

\emph{No participante}: la preferencia del inversor es \money{5000000}; la conversión pro
rata rendiría \(25\% \times \money{8000000} = \money{2000000}\). Aplicando
\cref{eq:liquidation-preference}:
\[
  \mathrm{Inversor}_{\text{no-part}}
  \;=\; \max(\money{5000000},\;\money{2000000})
  \;=\; \money{5000000}.
\]
Los fundadores y accionistas comunes reciben \(\money{8000000} - \money{5000000} = \money{3000000}\).

\emph{Participante}: el inversor primero toma \money{5000000} y luego participa en el
remanente:
\[
  \begin{aligned}
    \mathrm{Inversor}_{\text{part}}
    &\;=\; \money{5000000} \;+\; 0.25 \times (\money{8000000} - \money{5000000}) \\
    &\;=\; \money{5000000} \;+\; \money{750000} \;=\; \money{5750000}.
  \end{aligned}
\]
Los fundadores reciben \(\money{8000000} - \money{5750000} = \money{2250000}\). La cláusula
de participación transfirió \money{750000} de los fundadores al inversor --- el «doble cobro»
es de \money{750000} en este escenario, o aproximadamente el \qty{25}{\percent} del upside
incremental por encima del umbral de preferencia.

La magnitud del doble cobro crece con el tamaño de la salida: ante una salida de
\money{20000000}, el inversor participante tomaría \(\money{5000000} + 25\%\times\money{15000000} =
\money{8750000}\), frente a \(\max(\money{5000000},\,\money{5000000}) = \money{5000000}\)
no participante. La participación siempre tiene un costo para los fundadores en salidas
intermedias.
\end{example}

\begin{admonition}{Advertencia}
Las disposiciones antiditución componen la estructura de preferencia. La antidilución de
promedio ponderado de base amplia --- el estándar del mercado --- ajusta el precio de
conversión de las acciones preferentes existentes proporcionalmente si ocurre una ronda a
la baja, limitando, pero sin eliminar, el reprecio efectivo. La antidilución de trinquete
pleno reprecia el \emph{bloque completo} de preferentes previo al nuevo precio bajo
independientemente del volumen, lo que puede aniquilar el equity del fundador en una sola
ronda a la baja. La aritmética punitiva hace que las disposiciones de trinquete pleno sean
una señal de alerta desfavorable para el fundador; insista en el promedio ponderado de base
amplia \parencite{feld2019venture}. En mercados en dificultades, la combinación de una
preferencia participante \(1\times\) y antidilución de trinquete pleno puede dejar a los
accionistas comunes con ingresos casi nulos incluso a valores de salida moderados.
\end{admonition}

El análisis empírico de Kaplan y Strömberg sobre contratos de capital de riesgo muestra que
la asignación de derechos sobre flujos de efectivo (preferencias, antidilución), derechos de
control (asientos en la junta, vetos de votación) y derechos de liquidación se desacoplan
deliberadamente para alinear incentivos bajo incertidumbre \parencite{kaplan2003financial}.
Una preferencia que parece protectora al cierre de la ronda puede volverse punitiva al
momento de la salida; modelar la cascada completa a través de escenarios de salida plausibles
antes de firmar no es opcional.

\section{Métricas del Inversor y la Narrativa de Financiamiento}\label{sec:investor-metrics}
% complexity: 4

¿Cómo debe un fundador presentar el valor de la empresa a un inversor? La respuesta no es
sólo la narrativa; es la narrativa \emph{anclada en los números} que los inversores usan
realmente para comparar y fijar el precio de las startups. En 2026 el filtro cuantitativo
primario no es la tasa de crecimiento en aislamiento, sino la relación entre crecimiento y
consumo de capital, expresada a través de la Regla del 40 (\cref{eq:rule40}) y el múltiplo
de quema (\cref{eq:burn-multiple}). El ancla de valoración, a su vez, es el múltiplo de
EV/ARR de \cref{eq:ev-revenue}:
\begin{equation}\label{eq:arr-multiple}
  \mathrm{PreMoney} \;\approx\; \mathrm{ARR} \;\times\; \frac{\mathrm{EV}}{\mathrm{ARR}},
\end{equation}
donde el múltiplo de EV/ARR (\cref{eq:ev-revenue}) es el precio de equilibrio actual del
mercado por cada dólar de ingreso recurrente, ajustado por el perfil de eficiencia de
crecimiento de la empresa. Una empresa de alto crecimiento y eficiente en capital exige un
múltiplo premium; una empresa de crecimiento lento y alto burn enfrenta un descuento
pronunciado. La puntuación de la Regla del 40 y el múltiplo de quema son las dos palancas
que los inversores usan para mover el múltiplo \parencite{metrick2021venture}.

\begin{figure}[htb]
  \centering
  \includegraphics[width=0.92\linewidth]{cap-table-waterfall}
  \caption{%
    \textbf{Dilución de cap table y cascada de liquidación.}
    \emph{Panel izquierdo:} propiedad del fundador a través de tres rondas de financiamiento
    (Pre-seed \qty{100}{\percent} $\to$ Seed \qty{75}{\percent} $\to$ Series~A \qty{62.5}{\percent}).
    \emph{Panel derecho:} ingresos de salida para fundadores frente al inversor de Series~A
    bajo preferencias de liquidación \(1\times\) no participante y participante, para valores
    de salida de \money{5000000} a \money{30000000} (Series~A: \money{5000000} invertidos,
    \qty{25}{\percent} de propiedad). El punto de inflexión de la participación es visible
    alrededor de \money{20000000}, donde el valor de conversión del inversor supera la
    preferencia.%
  }
  \label{fig:cap-table-waterfall}
\end{figure}

\begin{example}[Narrativa de métricas del inversor para un pitch de Series~A]
La empresa de referencia: ARR = \money{4000000}; crecimiento de ARR en los últimos doce
meses = \qty{30}{\percent}; margen bruto = \qty{70}{\percent}. En aislamiento, un
crecimiento del \qty{30}{\percent} a este ARR podría sustentar un múltiplo de EV/ARR de
\(8\times\), lo que implica un pre-money de
\[
  \money{4000000} \;\times\; 8 \;=\; \money{32000000}.
\]
Pero el inversor examina \cref{eq:rule40}: puntuación de la Regla del 40 =
\qty{30}{\percent} de crecimiento $+ (-\qty{15}{\percent})$ de margen EBITDA $= 15$. Una
puntuación de 15 está muy por debajo del umbral de 40 puntos donde crecimiento y
rentabilidad están equilibrados; el inversor descuenta el múltiplo de $8\times$ a $6\times$:
\[
  \mathrm{PreMoney} \;=\; \money{4000000} \;\times\; 6 \;=\; \money{24000000}.
\]
El múltiplo de quema (\cref{eq:burn-multiple}): la empresa añadió \money{1200000} de ARR
nuevo neto (30\% de \money{4000000}) mientras quemaba \money{1500000} de efectivo neto, lo
que arroja un múltiplo de quema de \(\money{1500000}/\money{1200000} = 1.25\times\). A
$1.25\times$, el múltiplo de quema es aceptable --- por debajo del umbral de $2\times$ que
desencadena descuentos materiales --- y el fundador debería liderarlo. Combinado con el
pronóstico de \emph{runway} de \cref{eq:runway} (efectivo / tasa de quema neta), la
narrativa se convierte en: «Tenemos 18 meses de \emph{runway} al ritmo de quema actual, un
múltiplo de quema de $1.25\times$, y un camino claro hacia la paridad con la Regla del 40
en el mes 24 conforme se materialice el apalancamiento del margen bruto.»

El pre-money de \money{24000000} se sitúa cómodamente entre los escenarios base y optimista
de DCF de \cref{eq:scenario-ev} (\money{15000000}--\money{28000000}), proporcionando a los
inversores un ancla comparable al mercado para el DCF que ejecutarán de forma independiente.
Vincular el múltiplo al DCF elimina la apariencia de fijación arbitraria de precios y
centra la conversación en los fundamentos.
\end{example}

\begin{admonition}{Nota}
Los múltiplos de ARR varían significativamente según sector, nivel de crecimiento y
condiciones de mercado; la cifra de $6\times$ del ejemplo refleja un perfil de eficiencia
moderada y crecimiento medio en el entorno de 2026. Las empresas líderes que superan 60
puntos en la Regla del 40 han alcanzado múltiplos de \(10\times\)--\(15\times\) ARR. El
múltiplo apropiado no es el resultado de una fórmula, sino una observación de mercado; el
marco de EV/ARR de \cref{eq:ev-revenue} provee el ancla teórica (la tasa de crecimiento, el
margen y la tasa de descuento determinan el valor), mientras que las transacciones
comparables proporcionan la verificación empírica.
\end{admonition}

La narrativa de financiamiento es una capa de traducción: convierte la lógica DCF de
\cref{eq:dcf} en las métricas operativas que los inversores realmente rastrean. La Regla del
40, el múltiplo de quema y el múltiplo de ARR no son alternativas al pensamiento de flujo de
caja descontado; son atajos hacia la misma conclusión de valor presente. Un fundador que
entiende por qué \cref{eq:arr-multiple} es simplemente una versión comprimida de
\cref{eq:ev-revenue} --- que a su vez es una versión comprimida de \cref{eq:dcf} --- puede
defender cualquier métrica que cuestione un inversor sofisticado, porque la cadena conduce
de vuelta a los primeros principios. \textcite{metrick2021venture} documentan cómo los
inversores institucionales construyen sus propios modelos DCF junto con los múltiplos,
usando la brecha entre ambos como señal de si el fundador entiende su propio negocio.
